% chapters/06/section-02.tex
\section{Gravitasi dan Batas Lantai}
\label{sec:gravitasi-lantai}

Benda yang jatuh perlu berhenti di suatu tempat. Bagian ini membahas lantai
sebagai batas, dan persoalan yang muncul ketika benda menyentuhnya.

\subsection{Lantai sebagai Batas Bawah}

Lantai dinyatakan sebagai satu nilai tetap: jarak tepi atas lantai dari tepi
atas jendela. Pada program bab ini, lantai diletakkan enam puluh piksel di
atas tepi bawah jendela, menyisakan ruang untuk menampilkan keterangan.

Setiap bingkai, sesudah posisi diperbarui, program memeriksa apakah benda
sudah menyentuh atau melewati lantai.

\begin{kodekeluaran}
if position.y + half >= floor_y:
    position.y = floor_y - half        # dikembalikan ke permukaan lantai
    if velocity.y > 0:                 # sedang bergerak ke bawah
        velocity.y = -velocity.y * restitution
        if abs(velocity.y) < MIN_BOUNCE_SPEED:
            velocity.y = 0.0           # pantulan dihentikan
    hit_floor = True
\end{kodekeluaran}

Tiga hal perlu diperhatikan pada potongan kode tersebut.

\begin{enumerate}
  \item \textbf{Pemeriksaan memakai tanda ketidaksamaan, bukan kesamaan.}
        Menuliskan \code{if position.y + half == floor_y} akan gagal begitu
        benda bergerak cukup cepat, karena pada satu bingkai posisinya masih di
        atas lantai dan pada bingkai berikutnya sudah jauh di bawahnya. Benda
        itu akan menembus lantai tanpa pernah terdeteksi. Inilah gejala
        \istilah{tunneling} yang telah disebut pada Bagian~2.4, dan
        penanganannya secara lengkap dibahas pada Bab~VII.
  \item \textbf{Posisi dikembalikan ke permukaan lantai} sebelum kecepatan
        diubah. Tanpa langkah itu, benda akan terbenam sebagian ke dalam
        lantai dan lambat laun masuk lebih dalam.
  \item \textbf{Kecepatan ke bawah dinolkan atau dibalik}, bergantung pada
        koefisien restitusi. Bagian berikutnya membahas pilihan itu.
\end{enumerate}

\subsection{Mengapa Pemain Tidak Boleh Memantul}

Pada program bab ini, pemain dan bola memakai fungsi pembaruan yang sama,
tetapi satu nilai argumennya berbeda. Perhatikan akibatnya.

\begin{table}[htbp]
\centering
\caption{Akibat koefisien restitusi pada perilaku benda yang mendarat}
\label{tab:restitusi-benda}
{\setlength{\tabcolsep}{4pt}\footnotesize%
\begin{tabularx}{\textwidth}{@{}C{2.4cm} C{1.6cm} L{4.2cm} Y@{}}
\toprule
\headrow \textbf{Benda} & \textbf{Restitusi} & \textbf{Perilaku saat mendarat} & \textbf{Alasan} \\
\midrule
Pemain & $0{,}0$ &
Berhenti seketika, kecepatan tegaknya menjadi nol. &
Kendali pemain menuntut ketepatan. Pemain yang memantul setelah mendarat tidak
dapat mengarahkan geraknya, dan permainan terasa tidak dapat diandalkan. \\
\addlinespace
Bola & $0{,}8$ &
Memantul kembali dengan kecepatan 80 persen dari kecepatan datangnya. &
Pantulan adalah perilaku yang memang dikehendaki, dan kehilangan kecepatan
membuat benda akhirnya berhenti dengan sendirinya. \\
\bottomrule
\end{tabularx}}
\end{table}

Tidak ada kode tambahan untuk membedakan keduanya. Satu
parameter sudah cukup. Inilah manfaat menuliskan hukum fisika sebagai fungsi
yang menerima nilai sebagai argumen, sebagaimana disebutkan pada
Bagian~\ref{sec:percepatan}.

\begin{peringatanbox}
Kesalahan yang sering terjadi pada tahap ini adalah memberikan koefisien
restitusi kepada pemain ``supaya terlihat hidup''. Akibatnya pemain memantul
beberapa kali setiap kali mendarat, dan kendalinya terasa lepas. Apabila Anda
menghendaki kesan hidup pada gerak pemain, tambahkan percepatan dan
perlambatan pada gerak mendatarnya, bukan pantulan pada gerak tegaknya.
\end{peringatanbox}

\subsection{Mengenali Benda yang Menyentuh Lantai}

Lompatan pada bagian berikutnya hanya boleh dilakukan ketika benda berada di
tanah. Karena itu, keadaan itu perlu dikenali dengan pasti.

Cara yang dipakai pada program bab ini adalah memeriksa dua syarat sekaligus:
badan benda bersandar pada permukaan lantai, dan kecepatan tegaknya tidak lagi
mengarah ke bawah.

\begin{kodekeluaran}
on_ground = (position.y + half >= floor_y - 0.5 and velocity.y >= 0)
\end{kodekeluaran}

Toleransi setengah piksel diberikan karena posisi benda berupa bilangan
pecahan. Tanpa toleransi itu, benda yang mendarat pada ketinggian 459,7
sementara lantai berada pada 460,0 akan dianggap belum menyentuh lantai, dan
lompatan tidak dapat dilakukan walaupun benda itu tampak sudah berdiri di
atasnya.

\begin{catatabox}
Pada benda yang jatuh dari ketinggian besar, kecepatannya dapat tumbuh
melampaui batas yang masuk akal dan membuat benda menembus lantai dalam satu
bingkai. Cara penanganan yang sederhana adalah membatasi kecepatan jatuh pada
nilai tertentu, misalnya 1200 piksel per detik, sebelum posisi diperbarui.
Nilai batas itu disebut kecepatan akhir (\istilah{terminal velocity}) dan
banyak dipakai pada permainan.
\end{catatabox}
